\section{The Hazard Framework: A Structurally Distinct Estimator}\label{sec:hazard}

This is the estimator on which the paper's central decomposition rests. The ABM of Section~\ref{sec:abm} does not approach the hazard framework's aggregate recovery. On Section~\ref{sec:robustness-hybrid}'s shared accounting basis the microsimulation form recovers 91.3\% at the headline off-window floor and 97.9\% at the in-sample calibration point. The cohort-hazard form corroborates at 112.4\% but is excluded from headline ranges (Section~\ref{sec:patha}). The ABM's production specification recovers 11.9\% on the frozen draw (fifty-seed mean 13.6\%). Section~\ref{sec:robustness-crossdesign}'s cross-design leg leaves open whether that contrast is modeling paradigm or data source. Even the ABM's highest-recovery behavioral extension, at 54.9\%, falls well short. I attribute that figure in Section~\ref{sec:abm-results} to a mechanical recalibration artifact. The uncorrected rational baseline's 71.1\% predates that section's scheduled-amortization, CUSIP-parse, and CPR-back-out corrections. It is not in the production estimate family. Both hazard paths and the ABM face the same aggregate benchmark. Only Path B's recovery is accurate on the benchmark-consistent accounting basis. And no estimator's monthly co-movement with the empirical path survives removal of the window's shared trend, Path B's included (Section~\ref{sec:pathb}). Path A's monthly path fit is no better than the ABM's. Sections~\ref{sec:patha} and~\ref{sec:hazard-interp} quantify that limitation, because it bears on each recovery figure's evidential weight.

\subsection{Motivation}\label{sec:hazard-motivation}

Section~\ref{sec:abm}'s pattern is difficult to reconcile with the hypothesis that the residual reflects a missing rational or behavioral household mechanism that a sufficiently rich utility function would eventually capture. It is more consistent with the hypothesis that the binding constraints on the empirical prepayment path lie at least partly outside the household decision function: in loan-level heterogeneity, servicer pipeline dynamics, and cohort-level path dependence that a fixed, representative-household population struggles to represent.

The hazard framework introduced in Section~\ref{sec:method-estimators} is built to test this hypothesis. It abandons the household utility function---no mobility-desire parameter, no loss-aversion coefficient, no utility-maximizing decision rule---and models loan-level survival directly, using proportional hazards (proportional above a baseline turnover floor; Section~\ref{sec:pathb}) estimated on real loan-level data; the closed-population survivor-selection failure of Section~\ref{sec:abm-results} motivates a cohort-level treatment of burnout, not a per-loan one. For comparability the framework retains the same \$764.7 billion empirical benchmark and the same QT window definition as the ABM. It comprises two paths: a stratum-month discrete-time hazard model estimated on aggregated loan-level data, and a loan-level microsimulation calibrated to literature elasticities with explicit competing risks. I use competing risks in the sense of \citet{deng2000}, who estimate prepayment and default as dependent competing hazards over unobserved borrower heterogeneity.

\begin{table}[H]
\centering
\caption{Compact specification of the two hazard paths. Every entry restates
content specified in Sections~\ref{sec:patha} and~\ref{sec:pathb} or
Appendix~\ref{app:params}.}
\label{tab:specbox}
\footnotesize
\begin{threeparttable}
\begin{tabular}{@{}p{2.2cm}p{5.9cm}p{5.9cm}@{}}
\toprule
 & Path A (estimated cohort hazard) & Path B (literature-calibrated microsimulation) \\
\midrule
Unit; sample & stratum-month; full Freddie 2017--2021 origination universe (296 four-way cells) & loan-month; stratified 75{,}000-loan draw from the same universe \\
Outcome; estimation & prepaid UPB with log-exposure offset; Poisson pseudo-maximum likelihood, log link \eqref{eq:pathA} & monthly prepayment hazard \eqref{eq:pathB}; nothing estimated on the benchmark; every parameter fixed a priori, though the dominant level-setter (the turnover floor) is calibrated in-window (Sections~\ref{sec:pathb}, \ref{sec:robustness-floor}; Appendix~\ref{app:params}) \\
Age profile & seasoning spline, knots $\{12,24,36,60,84,120\}$ months (seven columns) & 100~PSA ramp to 6\% CPR at month 30 \\
Rate-gap term & estimated: $+0.65$ standardized ($+0.29$ per 100 bp), spec v4 & imported from \citet{liebersohn2024} via \eqref{eq:beta1}; 5.5--7.7\% band \\
Stratification & 295 stratum fixed effects (vintage $\times$ coupon $\times$ FICO $\times$ LTV) & stratum-level burnout state $B_{s(i),t}$; standardized FICO/LTV ex-ante priors \\
Time effects & eleven calendar-month dummies (spec v4) and the dynamic friction index & none; competing default hazard \eqref{eq:default} and the Markov delinquency pipeline; the omitted housing-activity level term is priced rather than modelled, by the scaled-null variant of Section~\ref{sec:identification} (run \texttt{scaled\_null\_housing\_activity}) \\
Turnover floor & none (stratum fixed effects set baseline levels) & 4\% CPR floor, max form; competing-risks form tested (Section~\ref{sec:robustness-floor}) \\
Inference & stratum-cluster and temporal moving-block bootstraps (Table~\ref{tab:bootstrap}); no claim survives them---the rate-gap coefficient's sign is not distinguishable from zero under any bias-respecting construction (Section~\ref{sec:patha}) & stratified loan-level bootstrap (200 replicates) for sampling; calibration box for parameters (Section~\ref{sec:robustness-floor}) \\
Validation & temporal holdout, reported as a diagnostic only (Section~\ref{sec:patha}) & Fannie Mae external replication (marginal $+8.68$ points against Freddie's $+9.20$; the pre-committed envelope it clears is 11.06 points wide) and the synthetic companion at 106.0\% (Sections~\ref{sec:pathb-fannie}, \ref{sec:robustness-synthesis}) \\
Role & aggregate context only; does not corroborate the result; excluded from headline figures & headline estimator: at the off-window floor central 91.3\%, $\beta_1 = 0$ null 85.7\%, marginal $+5.6$ points; at the in-sample calibration point 97.9\%, 88.7\%, $+9.2$ points (all shared basis) \\
\bottomrule
\end{tabular}
\end{threeparttable}
\end{table}
